\section{The integral lowering map for every auxiliary character}
\label{atl-section}

We extend the class comparison to the nontrivial auxiliary
characters.  The relevant operation increases the rank by one,
then restricts a coefficient representation to the smaller rank.
Its integral normalization involves the conjugate of the selected
prime generator.  We give the local groups, the extension of
decorated lattices, the transport between class representatives,
and the descent of the resulting complex-valued functions.
The polynomial branching theorem is the separate result
Theorem~\ref{ib-standard-integral}; no automorphic decomposition
or period identity is needed for the arguments in this section.
The definitions compared here are those of the actual final
step for every auxiliary character in the source
\cite{alpo-uniform-source}.

\subsection{The selected primes and the character at the level}

Let $E/F$ be a CM quadratic extension, write
$\mathcal O=\mathcal O_E$, and write $x\mapsto\bar x$ for
its nontrivial automorphism.  A set $S$ of rational primes
contains the two type primes and all the primes $T$ used to
increase the rank.  For one such $T$ fix $\pi\in\mathcal O$
with
\begin{equation}\label{atl-prime-generator}
 \pi\bar\pi=T,\qquad
 (\pi)=\prod_{P\in\mathcal T}P,\qquad
 (\pi)+(\bar\pi)=\mathcal O.
\end{equation}
The primes in $\mathcal T$ are distinct.  The type primes
form another finite set $\mathcal G$ of primes of $E$, all
coprime to $T$.

The ideal equality in \eqref{atl-prime-generator} follows
from the actual prime data, rather than requiring a new
selection.  Namely, write $[E:\mathbb Q]=2f$.  The construction
has exactly $f$ distinct prime divisors of $\pi$, each of norm
$T$.  Conjugation preserves absolute ideal norms, so
$N((\pi))^2=N((T))=T^{2f}$ and $N((\pi))=T^f$.
The product of those $f$ prime divisors divides $(\pi)$
and has the same norm, hence equals $(\pi)$.  The coprimality
with its conjugate is also part of the selected prime data.

An auxiliary character determines a finite set $\mathcal Q$
of primes of $E$.  Each $Q\in\mathcal Q$ is the chosen prime
above a prime $v$ of $F$ such that
\begin{equation}\label{atl-good-primes}
 v\mathcal O=Q\bar Q,\quad Q\ne\bar Q,\quad
 N_{F/\mathbb Q}v=p\text{ is prime},\quad
 p\notin v^2,\quad p\notin S.
\end{equation}
The choice of one prime over each $v$ is fixed before the
auxiliary character is varied.  In particular no two selected
primes are conjugate.  The sets $\mathcal G$, $\mathcal T$,
and $\mathcal Q$ are disjoint, and $\mathcal Q$ is disjoint
from the conjugates of the other two sets.

\begin{lemma}\label{atl-prime-units}
For every $Q$ satisfying \eqref{atl-good-primes}, its residue
field has cardinality $p$, and $v_Q(p)=1$.  Both $\pi$ and
$\bar\pi$ are units in $\mathcal O_{(Q)}$.
These conclusions remain valid in residue characteristic two
whenever the displayed hypotheses hold.
\end{lemma}
\begin{proof}
The factorization of $v\mathcal O$ has ramification index one
at both its primes.  Since $[E:F]=2$, the equality
$\sum_{Q\mid v}e(Q/v)f(Q/v)=2$ gives residue degree one
at each prime.  Thus $NQ=Nv=p$.  Since $p$ belongs to $v$
and does not belong to $v^2$, its $v$-adic valuation is one;
extension to $Q$ leaves that valuation unchanged.

The rational prime $T$ differs from $p$, so $T$ is a unit
at $Q$.  Its two integral factors $\pi$ and $\bar\pi$
must then both have valuation zero.  This proves the last
assertion as well, without using the residue characteristic.
\end{proof}

Fix integers $m\geq2$, $r,t_0\geq0$, and put
$n=r+mt_0$.  At $Q\in\mathcal Q$, let
$\kappa_Q=\mathcal O/Q$ and let
$\xi_Q:\kappa_Q^\times\to K^\times$ be the residue
character, where the number field $K$ contains its values.
The lower-rank residue group and character are
\begin{equation}\label{atl-row-level}
 H_Q=\left\{
 \begin{pmatrix}A&B\\0&D\end{pmatrix}:
 A\in\operatorname{GL}_r(\kappa_Q),\quad
 D\in\operatorname{GL}_{mt_0}(\kappa_Q)\right\},
 \qquad
 \chi_Q\begin{pmatrix}A&B\\0&D\end{pmatrix}
                         =\xi_Q(\det D)^{-1}.
\end{equation}
An empty diagonal block has determinant one.  At a type prime
$P\in\mathcal G$ we keep its specified finite subgroup $H_P$
and character $\chi_P$.  The character of the action on
decorated objects is the product of these characters.  By
\eqref{acl-character}, the corresponding right character on
adelic functions is its inverse.  Thus the right character
at $Q$ is $\xi_Q(\det D)$, whereas the character in
\eqref{atl-row-level} is its reciprocal.

Here is the exact higher-rank local group.  Put
$A_Q=\mathcal O_{(Q)}$, let $\widetilde H_Q$ be the inverse
image of $H_Q$ in $\operatorname{GL}_n(A_Q)$, and set
\[
 K_{n+1}(C)=\ker\bigl(\operatorname{GL}_{n+1}(A_Q)
                 \longrightarrow\operatorname{GL}_{n+1}(A_Q/p^CA_Q)\bigr)
                    \qquad(C\geq1).
\]
Lemma~\ref{atl-prime-units} makes this the usual principal
congruence subgroup of depth $C$.

\begin{lemma}\label{atl-congruence-group}
The set
\begin{equation}\label{atl-top-group}
 \widetilde H_Q^{+}(C)
       =\operatorname{diag}(\widetilde H_Q,1)K_{n+1}(C)
\end{equation}
is a subgroup of finite index in $\operatorname{GL}_{n+1}(A_Q)$.
An element $k$ belongs to it if and only if there exists
$h\in\widetilde H_Q$ with
\[
                      k\equiv\operatorname{diag}(h,1)\pmod {p^C}.
\]
The formula
\begin{equation}\label{atl-top-character}
 \chi_Q^{+}(k)=\chi_Q(\bar h)
       \quad\text{if }k=\operatorname{diag}(h,1)u,
                         \quad u\in K_{n+1}(C),
\end{equation}
is a well-defined character.  It is trivial on $K_{n+1}(C)$,
and its restriction to $\operatorname{diag}(\widetilde H_Q,1)$
is exactly the lower-rank character.  Furthermore
\[
 \operatorname{diag}(h,1)\in\widetilde H_Q^{+}(C)
                       \quad\Longleftrightarrow\quad
 h\in\widetilde H_Q
 \qquad(h\in\operatorname{GL}_n(A_Q)).
\]
\end{lemma}
\begin{proof}
The principal congruence subgroup is normal, so its product
with a subgroup is a subgroup.  It has finite index because
the quotient ring $A_Q/p^CA_Q$ is finite.  Multiplication
by an integral invertible matrix preserves $p^CA_Q^{n+1}$,
so the membership criterion is equivalent to the displayed
factorization.

Reduction modulo $Q$ of every element of the product is
$\operatorname{diag}(\bar h,1)$ with $\bar h\in H_Q$.
It determines $\bar h$ uniquely.  Pulling $\chi_Q$ back
along this homomorphism proves well-definedness and all
the character assertions.  The same reduction proves the
forward implication in the last equivalence; the reverse
implication uses $u=1$.
\end{proof}

For the actual choice of depth, if $M\geq0$ is the local
constant used for this auxiliary character, then
\begin{equation}\label{atl-depth}
                  C=\max(4m+1,2m+2M).
\end{equation}
Indeed the preceding depth is $\max(4m,2m+2M-1)$, and
adding one gives \eqref{atl-depth}; $m\geq2$ ensures that
the subtraction in the second expression does not cause
an exceptional case.  The group in \eqref{atl-top-group}
is not replaced by its larger inverse image under reduction
modulo $Q$.  Its character is, however, already determined
by that reduction.  At the type primes the higher-rank
residue group is simply $\operatorname{diag}(H_P,1)$.

\subsection{Extension by the selected line}

Use the standard positive Hermitian form on $E^n$, with
the convention $h(v,w)=\bar v^{t}w$.  Let $N\subset E^n$
be a full self-dual $\mathcal O$-lattice.  A lower-rank
decorated object consists of $N$, its residue frames
\[
 b_P:N/PN\xrightarrow{\sim}\kappa_P^n
                       \quad(P\in\mathcal G),
\]
and at every $Q\in\mathcal Q$ a basis matrix
$F_Q\in\operatorname{GL}_n(E)$ with
$F_QA_Q^n=N_Q$.  Replacing all local bases by their
principal-congruence orbits gives the equivalent residue
description of Proposition~\ref{acl-row-classes}.

Extend this object to rank $n+1$ by
\begin{align}
 N^+&=N\oplus\mathcal O,\notag\\
 \ell_P(v,s)&=\bar\pi s\bmod P
                    &&(P\in\mathcal T),\label{atl-extended-object}\\
 b_P^+(v,s)&=(b_P(v),s/\pi\bmod P)
                    &&(P\in\mathcal G),\notag\\
 F_Q^+&=\operatorname{diag}(F_Q,\pi)
                    &&(Q\in\mathcal Q).\notag
\end{align}
Here $\ell_P$ is an $\mathcal O$-linear map to $\kappa_P$.
All divisions in the frames are by local units: at the
type primes this follows from $T$ being different from
the type primes, and at $Q$ it follows from
Lemma~\ref{atl-prime-units}.

\begin{proposition}\label{atl-extension}
The formulas \eqref{atl-extended-object} define a decorated
unimodular lattice.  Its lattice cut out by the functionals,
and its Hermitian dual, are respectively
\begin{equation}\label{atl-kernel-dual}
 L^+=N\oplus\pi\mathcal O,
               \qquad (L^+)^\#=N\oplus\bar\pi^{-1}\mathcal O.
\end{equation}
The construction is injective on decorated objects and
equivariant under
\[
 \gamma\longmapsto\operatorname{diag}(\gamma,1),\qquad
 h_P\longmapsto\operatorname{diag}(h_P,1),\qquad
 k_Q\longmapsto\operatorname{diag}(k_Q,1).
\]
These homomorphisms are injective, preserve the level
characters exactly, and send every lower-rank stabilizer
into the stabilizer of the extended object.  In the
parameter used for branching one has
\begin{equation}\label{atl-tz}
                            t=\bar\pi,\qquad \bar t t=T.
\end{equation}
\end{proposition}
\begin{proof}
The direct sum of two self-dual lattices is self-dual.
At $P\mid\pi$, the element $\bar\pi$ is a unit by
\eqref{atl-prime-generator}, so $\ell_P$ is nonzero.
Its kernel imposes precisely $s\in P$.  Intersecting
these conditions gives $s\in\prod_{P\in\mathcal T}P=(\pi)$.
This proves the first formula in \eqref{atl-kernel-dual}.
The pairing of $(0,\pi a)$ with $(v,s)$ is
$\bar\pi\bar a s$, so the last coordinate of the dual
is $\bar\pi^{-1}\mathcal O$; its first coordinates
form $N^\#=N$.

The displayed frames are invertible because their last
diagonal entries are local units.  The original lattice,
residue frames, and local bases are recovered by taking
the first $n$ coordinates, proving injectivity on objects.

For the residue-frame action
$b_P\mapsto h_Pb_P\gamma^{-1}$, its extension is
$b_P^+\mapsto\operatorname{diag}(h_P,1)b_P^+
\operatorname{diag}(\gamma,1)^{-1}$.
For a local basis the action is $F_Q\mapsto\gamma F_Qk_Q^{-1}$;
therefore
\[
 \operatorname{diag}(\gamma,1)\operatorname{diag}(F_Q,\pi)
       \operatorname{diag}(k_Q,1)^{-1}
                    =\operatorname{diag}(\gamma F_Qk_Q^{-1},\pi).
\]
The functionals at $\mathcal T$ are unchanged by the
unitary embedding.  These identities prove equivariance.
Injectivity of the group embeddings is immediate from
their first blocks.  Preservation of characters is
Lemma~\ref{atl-congruence-group} at $Q$, and the definition
of the diagonal residue character at the type primes.
Equivariance then proves the stabilizer assertion.
Finally \eqref{atl-tz} is exactly the last coordinate
of the dual lattice, together with $\pi\bar\pi=T$.
\end{proof}

This proposition does not assert that the induced map of
class sets is injective or surjective.  Both properties
are unnecessary.  At any class which meets the image one
may choose an extended object as representative; all
other classes keep arbitrary representatives.  There
is consequently no need to assume that every higher-rank
object comes from the smaller rank.

Both class sets at these levels are finite.  Indeed, the
finiteness theorem for positive definite unimodular lattices
gives finitely many underlying lattices up to rational
unitary isometry.  For each one, every residue functional
and residue frame has only finitely many possibilities.
The set of bases at $Q$ is a torsor under the integral
linear group.  Its quotient by the chosen level subgroup
is finite by Lemma~\ref{atl-congruence-group}, and likewise
for the lower-rank inverse image of $H_Q$.  Taking the
product over the finite set of primes leaves finitely
many possibilities.  This argument also applies to
higher-rank classes outside the image of the extension.

\subsection{Simultaneous bases and the scaled branching map}

Retain the coefficient field $K$, the finite family of
embeddings $\tau:E\hookrightarrow K$, and the representations
$R_\lambda$ and integral coefficient lattices of
\eqref{acl-contragredient}.  For a full lattice $M\subset E^d$
abbreviate its contragredient coefficient lattice by
$\Delta_\lambda(M)$.  It is the integral coefficient lattice used in the
class construction; all determinant twists are included.

\begin{lemma}\label{atl-semilocal-basis}
Let $\ell$ be a rational prime and $M\subset E^d$ a full
$\mathcal O$-lattice.  There exists
$B\in\operatorname{GL}_d(E)$ whose columns are a basis of
$M_P$ over $\mathcal O_{(P)}$ simultaneously at every
$P\mid\ell$.  If $w\mid\ell$ is a coefficient prime of $K$,
then
\begin{equation}\label{atl-local-coefficient}
 v\in\Delta_\lambda(M)_w
 \quad\Longleftrightarrow\quad
 R_\lambda(B^{-1})v\in
                 \Delta_\lambda^{\rm std}(\mathcal O_{K,w}).
\end{equation}
\end{lemma}
\begin{proof}
For every $P\mid\ell$, the module $M_P$ is free of rank
$d$ over the discrete valuation ring $\mathcal O_{(P)}$.
Choose a basis of $M/PM$.  The Chinese remainder theorem
for the pairwise comaximal submodules $PM$ lets us choose
$d$ elements of $M$ with these prescribed residue bases
at all the finitely many primes above $\ell$.
At each such prime their span has quotient $C$ with
$C/PC=0$.  The quotient is finitely generated; Nakayama's
lemma gives $C=0$.  Since $M_P$ has rank $d$, the selected
elements are a basis.  They are linearly independent over
$E$, and hence form the columns of an invertible matrix
$B$.  For $d=0$ take the unique empty matrix.

For each embedding $\tau$ the prime $\tau^{-1}(w)$ is
above $\ell$, so this same matrix supplies every required
local basis.  The defining change-of-basis rule from
Section~\ref{coef:section} gives
\[
 \Delta_\lambda(M)_w
     =R_\lambda(B)\Delta_\lambda^{\rm std}(\mathcal O_{K,w}).
\]
Applying $R_\lambda(B^{-1})$ proves the equivalence.
\end{proof}

Let $\lambda_\tau$ and $\sigma_\tau$ be dominant integer
weights in ranks $n+1$ and $n$, respectively.  We require
the actual interlacing and determinant-shift conditions
\begin{equation}\label{atl-weights}
 \lambda_{\tau,i}\geq\sigma_{\tau,i}
                       \geq\lambda_{\tau,i+1}\ (1\leq i\leq n),
               \qquad \sigma_{\tau,n}=\lambda_{\tau,n+1}.
\end{equation}
The ranks in the application are positive.  The standard
branching theorem gives a $K$-linear map
\[
 B_1:W_\lambda(K)\longrightarrow W_\sigma(K)
\]
with the specified polynomial normalization, such that
\begin{align}
 B_1R_\lambda(\operatorname{diag}(\gamma,1))
                &=R_\sigma(\gamma)B_1,
                    &&\gamma\in\operatorname{GL}_n(E),
                         \label{atl-standard-covariance}\\
 B_1\Delta_\lambda^{\rm std}(\mathcal O_{K,w})
                &\subseteq\Delta_\sigma^{\rm std}(\mathcal O_{K,w})
                    &&\text{for every }w.
                         \label{atl-standard-integral}
\end{align}
These are precisely the integral assertion of
Theorem~\ref{ib-standard-integral} and the diagonal
covariance of Corollary~\ref{ib-diagonal-covariance}, applied
to the contragredient realization.  Define
\begin{equation}\label{atl-scaled-map}
                 B_t=B_1R_\lambda(\operatorname{diag}(I_n,t))
                         \qquad(t\in E^\times).
\end{equation}
To check the normalization explicitly, set
\[
 b_{\tau,i}=\lambda_{\tau,i}-\sigma_{\tau,i}\geq0,
 \qquad a_\tau=\lambda_{\tau,n+1},\qquad
 N_\tau=\sum_{i=1}^{n+1}\lambda_{\tau,i}
                              -\sum_{i=1}^n\sigma_{\tau,i}.
\]
The standard polynomial formula first retains exactly the
monomials with exponent $b_{\tau,i}$ at the last-row variable
$X_{\tau,n+1,i}$, and exponent zero at every last-column
variable; it then applies a fixed projection and scalar.
Under $R_\lambda(\operatorname{diag}(I_n,t))$, each retained
monomial is multiplied by
$\prod_\tau\tau(t)^{-a_\tau-\sum_i b_{\tau,i}}$.
Since $a_\tau+\sum_i b_{\tau,i}=N_\tau$, this proves the
exact formula
\begin{equation}\label{atl-scaled-factor}
                      B_t=\left(\prod_\tau\tau(t)^{-N_\tau}\right)B_1.
\end{equation}
The exponents $N_\tau$ are integers and need not be
nonnegative.  This is the scaled map in the construction,
including the determinant contribution $a_\tau$.

\begin{proposition}\label{atl-scaled-integrality}
For $B'\in\operatorname{GL}_n(E)$ put
$B=\operatorname{diag}(B',t^{-1})$.  Then
\begin{equation}\label{atl-coordinate-identity}
           R_\sigma(B'^{-1})B_t=B_1R_\lambda(B^{-1}).
\end{equation}
The map $B_t$ intertwines the two representations under
$\gamma\mapsto\operatorname{diag}(\gamma,1)$.
For every full lattice $M\subset E^n$ it satisfies
\begin{equation}\label{atl-branch-integral}
 B_t\Delta_\lambda(M\oplus t^{-1}\mathcal O)
                           \subseteq\Delta_\sigma(M).
\end{equation}
All these statements hold at every finite coefficient
prime, including primes above $T$, primes above the
auxiliary set $\mathcal Q$, and primes above two.
\end{proposition}
\begin{proof}
By \eqref{atl-standard-covariance}, the left side of
\eqref{atl-coordinate-identity} is
\[
 B_1R_\lambda(\operatorname{diag}(B'^{-1},1))
       R_\lambda(\operatorname{diag}(I_n,t))
                 =B_1R_\lambda(\operatorname{diag}(B'^{-1},t)).
\]
This is the right side.  The same covariance, and the
commutation of the last-coordinate scaling with
$\operatorname{diag}(\gamma,1)$, prove the intertwining
assertion.

Fix a coefficient prime $w\mid\ell$.  Choose $B'$ as
in Lemma~\ref{atl-semilocal-basis} for $M$.  The matrix
$B=\operatorname{diag}(B',t^{-1})$ is a simultaneous
basis of $M\oplus t^{-1}\mathcal O$.  By
\eqref{atl-local-coefficient}, a vector in the localized
source lattice is carried by $R_\lambda(B^{-1})$ into
the standard integral lattice.  Apply
\eqref{atl-standard-integral} and
\eqref{atl-coordinate-identity}, and use the same local
criterion for the target.  This proves the localized
inclusion at $w$.  Intersecting over all $w$ proves
\eqref{atl-branch-integral}.  The argument never requires
that $t$ be a unit at $w$, or that the residue
characteristic exceed the degree of a polynomial.
\end{proof}

Adding an integer $a_\tau$ to every entry of a weight
does not change its normalized partition
$\lambda_{\tau,i}-\lambda_{\tau,n+1}$, and therefore
does not change its standard polynomial lattice as a
set of polynomials.  It does change the determinant
factor in $R_\lambda$.  In \eqref{atl-coordinate-identity}
the representations are always those of the actual
shifted weights; only the standard polynomial lattices
may be identified by equality of normalized partitions.

\subsection{Transport of class values and global descent}

Let $\mathcal D$ be the lower-rank decorated objects and
$\mathcal D^+$ the higher-rank objects, with their groups
$G$ and $G^+$ acting as above.  Write
$j:\mathcal D\to\mathcal D^+$ for the extension and
$\iota:G\to G^+$ for the group homomorphism.  Let
$\chi,\chi^+$ be the action characters.  We have proved
\begin{equation}\label{atl-action-compatibility}
 j(gd)=\iota(g)j(d),\qquad
 \chi^+(\iota(g))=\chi(g),\qquad
 \gamma(\iota(g))=\operatorname{diag}(\gamma(g),1).
\end{equation}
For a higher-rank equivariant coefficient function $F$,
define
\begin{equation}\label{atl-lowering}
                         (\mathcal L_tF)(d)=B_tF(j(d)).
\end{equation}

\begin{theorem}\label{atl-global-lowering}
The map $\mathcal L_t$ carries higher-rank equivariant
$K$-valued forms to lower-rank equivariant $K$-valued
forms.  With $t=\bar\pi$ and the actual coefficient
lattices, it carries the integral class lattice into
the integral lower-rank class lattice, and the same
holds after localization at every coefficient prime.
It agrees after complexification with evaluation on
the extended objects followed by the normalized
polynomial branching map.

More explicitly, choose representatives $d_c$ of the
lower-rank classes and $e_b$ of the higher-rank classes.
Let $b(c)$ be the class of $j(d_c)$, and choose
$a_c\in G^+$ with $a_ce_{b(c)}=j(d_c)$.  If $v_b=F(e_b)$,
then the value at $d_c$ is exactly
\begin{equation}\label{atl-representative-formula}
             (\mathcal L_tF)(d_c)
                 =B_t\bigl(\chi^+(a_c)
                          R_\lambda(\gamma(a_c))v_{b(c)}\bigr).
\end{equation}
It is independent of the choices $a_c$.
\end{theorem}
\begin{proof}
Equations \eqref{atl-action-compatibility} and the
intertwining assertion of
Proposition~\ref{atl-scaled-integrality} give
\begin{align*}
 (\mathcal L_tF)(gd)
  &=B_t\chi^+(\iota(g))
               R_\lambda(\operatorname{diag}(\gamma(g),1))F(j(d))\\
  &=\chi(g)R_\sigma(\gamma(g))(\mathcal L_tF)(d).
\end{align*}
Thus the function has the required equivariance.

At $j(d)$ the coefficient lattice is that of
$(L_{j(d)})^\#=N_d\oplus t^{-1}\mathcal O$,
whereas the lower-rank lattice is that of $N_d$.
Their integral compatibility is
\eqref{atl-branch-integral}.  Transport from $e_{b(c)}$
to $j(d_c)$ preserves these exact lattices: the
unitary action carries the underlying lattice to
the required one, its coefficient action is $R_\lambda$,
and $\chi^+(a_c)$ is a root of unity, hence an integral
unit at every coefficient prime.  This proves the
global and localized assertions.

If $a_c'$ is another choice, $a_c^{-1}a_c'$ stabilizes
$e_{b(c)}$.  Its coefficient action fixes $v_{b(c)}$
by Proposition~\ref{acl-class-evaluation}, so the
transported values are equal.  The same argument,
or the equivariance just proved, shows that the
right side of \eqref{atl-representative-formula}
is fixed by the lower-rank stabilizer.  No averaging
over a stabilizer occurs.  Finally all coefficient
operations and \eqref{acl-evaluation} commute with
the chosen embedding $K\hookrightarrow\mathbb C$.
This proves the complex comparison.
\end{proof}

For completeness, the integral coordinates used to
recognize the resulting complex form can also be
characterized without presupposing rationality.

\begin{proposition}\label{atl-complex-descent}
Fix $w\mid\ell$ in $K$ and an embedding
$\jmath:K\hookrightarrow\mathbb C$.  At every lower-rank
class representative $d_c$, choose a simultaneous
basis $B_c'$ of $N_{d_c}$ as in
Lemma~\ref{atl-semilocal-basis}.  Let $z$ be a complex
equivariant form in the coefficient representation.
Extract its unique polynomial coefficient $P_c$ from
$z(d_c)$ by \eqref{acl-evaluation}.  Then $z$ is the
complexification of an element of the localized integral
class lattice if and only if, for every $c$,
\begin{equation}\label{atl-complex-criterion}
 R_\sigma^{\mathbb C}(B_c'^{-1})P_c
      \in\jmath\bigl(\Delta_\sigma^{\rm std}
                                        (\mathcal O_{K,w})\bigr).
\end{equation}
The complex representation is obtained through the
embeddings $\jmath\tau$.
\end{proposition}
\begin{proof}
If the form is already the complexification of a
localized integral one, the criterion is
\eqref{atl-local-coefficient} at every class.
Conversely, the right side of
\eqref{atl-complex-criterion} supplies a unique
$K$-valued polynomial $q_c$ in the standard local
lattice.  The uniqueness follows from injectivity
of the field embedding on polynomial coefficients.
Undoing $R_\sigma(B_c'^{-1})$ gives
$v_c=R_\sigma(B_c')q_c\in\Delta_\sigma(N_{d_c})_w$
whose complex evaluation is $z(d_c)$.

For any element of the stabilizer of $d_c$, the
complex evaluations of $v_c$ and its coefficient
translate are equal, since $z$ is equivariant.
The injectivity proved after \eqref{acl-evaluation}
implies equality over $K$.  Hence $v_c$ is fixed
by the entire stabilizer.  These values determine
a $K$-valued equivariant form by
Proposition~\ref{acl-class-evaluation}.

There are finitely many classes.  For each $c$ choose
$s_c\in\mathcal O_K\setminus w$ with
$s_cv_c\in\Delta_\sigma(N_{d_c})$.  The product
$s=\prod_cs_c$ is still outside $w$ and makes every
class value integral.  Multiplication by $s$ preserves
all stabilizer equalities.  Thus the form belongs to
the localized integral class lattice, and its
complexification is $z$.  This also proves uniqueness.
\end{proof}

To see the exact coordinate formula for lowering,
put $B_c=\operatorname{diag}(B_c',t^{-1})$.  For the
transported value in \eqref{atl-representative-formula}
define
\begin{align*}
 \operatorname{co}_+(F,c)
  &=\jmath\left(R_\lambda(B_c^{-1})
       \chi^+(a_c)R_\lambda(\gamma(a_c))v_{b(c)}\right),\\
 \operatorname{co}_-(z,c)
  &=R_\sigma^{\mathbb C}(B_c'^{-1})P_c.
\end{align*}
Equation \eqref{atl-coordinate-identity} gives
\begin{equation}\label{atl-final-coordinate}
 \operatorname{co}_-(\mathcal L_tF,c)
                   =B_1^{\mathbb C}\operatorname{co}_+(F,c).
\end{equation}
The first coordinate map carries every localized
integral form into the standard integral polynomial
lattice.  The second recognizes precisely the
localized integral forms by
Proposition~\ref{atl-complex-descent}.  This is the
full coordinate comparison needed in the final
lowering step.

\subsection{Uniformity over the required characters and scope}

Fix the finite family of auxiliary characters used
in the construction.  For each member use its actual
finite support $\mathcal Q$ and its own constant $M$
in \eqref{atl-depth}.  The preceding proofs apply
to every member, and to every selected prime $T$,
with exactly the same weights for fixed $T$-index.
The support and depth change the finite class set and
its level subgroup; they do not change the formulas
\eqref{atl-kernel-dual}, \eqref{atl-tz}, or the
standard polynomial map $B_1$.  The simultaneous
local bases may depend on the class, the character,
and $w$.  Their existence is proved after these
parameters are fixed and requires no new prime
selection or conductor estimate.

In particular, no exclusion $w\notin\mathcal Q$ is
needed for the integral lowering map itself.  Any
such exclusion in a later period-valuation theorem
is a hypothesis of that theorem, not a consequence
of this coordinate comparison.  The exact integral
polynomial branching assertion remains the dependency
\eqref{atl-standard-integral}; the preceding proofs
do not infer it merely from irreducibility over $K$.
This section does not prove the transfer identities
for every intermediate nontrivially twisted theta
sum, nor the valuation improvement in the period
pairing.  It does prove the actual final class map,
its characters, its local integral coordinates, and
their global descent.

\subsection{Checks on essential hypotheses}

The following examples explain why several exact
features cannot be omitted.  They do not contradict
the hypotheses of the construction.

First, at a prime of residue characteristic $p$,
take $n=1$, $H_Q=\operatorname{GL}_1(\kappa_Q)$, and
$C=2$.  The matrix $I+pE_{21}$ reduces to the identity,
so it lies in the full inverse image of the diagonal
residue group.  It does not lie in
$\operatorname{diag}(\widetilde H_Q,1)K_2(2)$, because
its lower left entry is not divisible by $p^2$.
Thus replacing the actual group by that inverse
image changes the level, even when the character
is unchanged.

Second, in $E=\mathbb Q(i)$ take $\pi=2+i$.
For the standard Hermitian form on one dimension,
$(\pi\mathbb Z[i])^\#=\bar\pi^{-1}\mathbb Z[i]$.
The proposed replacement $\pi^{-1}\mathbb Z[i]$
is different: $\bar\pi/\pi=(3-4i)/5$ is not an
algebraic integer.  The conjugation in $t=\bar\pi$
therefore affects integral lattices.

Third, transport between class representatives
cannot be omitted from
\eqref{atl-representative-formula}.  Already for a
free transitive action of the group $\{1,s\}$ on
two points, the equivariance rule $F(sd)=-F(d)$
gives opposite values at the two representatives
of the same orbit.  A branching map equal to the
identity would preserve this discrepancy.

Finally the finiteness of the class set is used in
the descent proof.  In an infinite product over
distinct rational primes $p_j\ne\ell$, the vector
$(1/p_j)_j$ belongs componentwise to $\mathbb Z_{(\ell)}$.
There is no integer $s$ prime to $\ell$ which makes
all its components integers.  Hence it does not
belong to $(\prod_j\mathbb Z)\otimes\mathbb Z_{(\ell)}$,
although it belongs to $\prod_j\mathbb Z_{(\ell)}$.
The finite class theorem and the common denominator
in Proposition~\ref{atl-complex-descent} prevent
this failure in the actual construction.


